\section{Related Work}

\subsection{Occupancy Forecasting}

Occupancy forecasting is a mature building-analytics problem rather than a new modeling task. Reviews by Jin \emph{et al.} and Caballero-Pe\~na \emph{et al.} organize the literature around forecasting applications, data sources, and evaluation processes~\cite{Jin2021,CaballeroPena2024}. Earlier commercial-building comparisons include stochastic and machine-learning approaches~\cite{LiDong2018}, while work based on indoor environmental measurements illustrates the potential and data-quality dependence of sensor-assisted prediction~\cite{RyuMoon2016}. Sequence models also remain active in the field, including Transformer-based occupancy prediction~\cite{Sun2023}. These studies motivate a broad baseline set, but they do not establish that one model family is uniformly preferable for every operating rule or building.

The present study accordingly retains a train-only schedule prior alongside nonlinear tabular and sequence models. The tabular reference uses LightGBM~\cite{Ke2017}, the random-forest comparator follows Breiman~\cite{Breiman2001}, and the sequence comparator is a compact encoder-only Transformer with a known-future calendar projection, architecturally motivated by~\cite{Vaswani2017}. The saved artifact identifier \texttt{DLinear} is a direct one-layer occupancy-history baseline, not the decomposition-based DLinear architecture of Zeng \emph{et al.}~\cite{Zeng2023}; it is retained only as a simple linear comparator and is relabelled accordingly in the paper. The goal is not to claim algorithmic novelty for any component. Instead, the comparison tests whether combining distinct saved score forecasts produces an inspectable, validation-selected offline policy under the same target and time split.

\subsection{From Occupancy Prediction to Building Operation}

Several studies couple occupancy predictions to a downstream HVAC or ventilation decision. Peng \emph{et al.} evaluate occupancy-prediction-based cooling control~\cite{Peng2018}. Esrafilian-Najafabadi and Haghighat analyze how occupancy-model choices relate to HVAC control performance~\cite{EsrafilianNajafabadi2022}, and Zhuang \emph{et al.} develop probabilistic occupancy forecasts for risk-aware ventilation~\cite{Zhuang2022}. Wietzke \emph{et al.} further connect occupancy prediction to model predictive control and simulation~\cite{Wietzke2024}. Those papers demonstrate why a forecast-only score may not summarize a downstream trade-off. They also delimit this paper's evidence: it does not execute an HVAC action, run a thermal simulation, or measure a counterfactual energy or comfort outcome.

The distinction is especially important when comparing offline studies with operational-system evaluations. Lin \emph{et al.} report an occupancy-based HVAC control-system evaluation in an office building~\cite{Lin2023}. This paper has no installation, occupant study, or measured before-and-after intervention. It therefore uses the term \emph{opportunity} only for a processed meter-derived load proxy coinciding with a recommended and camera-label-empty interval, and it reserves claims about savings or control performance for future intervention or simulation work.

\subsection{Decision-Layer Evaluation}

Decision-focused learning and predict-then-optimize methods train predictive models against a downstream decision objective~\cite{Wilder2019,ElmachtoubGrigas2022}. The present pipeline does not do this: model parameters are not trained with a differentiable decision loss. It fixes candidate score models, selects blend weights by validation AUPRC, and selects a score threshold by an empirical validation conflict cutoff and offline opportunity objective. The appropriate description is therefore a \emph{validation-selected decision-layer evaluation}, not a calibrated risk guarantee or an implementation of decision-focused learning.
